\documentclass[a4paper,12pt]{article}

\usepackage[T2A]{fontenc}
\usepackage[utf8]{inputenc}
\usepackage[russian]{babel}

\usepackage{amsmath,amssymb,amsthm,mathtools}
\usepackage{helvet}
\renewcommand{\familydefault}{\sfdefault}
\usepackage{icomma}
\usepackage{geometry}
\geometry{margin=2cm}
\usepackage{microtype}
\usepackage{tikz}
\usepackage{tkz-euclide}
\usepackage{pgfplots}
\pgfplotsset{compat=1.18}
\usepackage{tabularx,booktabs,longtable}
\usepackage{enumitem}
\usepackage{hyperref}
\usepackage{parskip}
\usepackage{fancyhdr}
\usepackage{setspace}
\usepackage{xcolor}
\usepackage{array}

\definecolor{accent}{HTML}{008080}
\definecolor{darkaccent}{HTML}{004D4D}
\definecolor{textgray}{HTML}{333333}
\definecolor{softgray}{HTML}{6E6E6E}

\hypersetup{
  colorlinks=true,
  linkcolor=darkaccent,
  urlcolor=darkaccent,
  citecolor=darkaccent
}

\onehalfspacing
\setlength{\parindent}{0pt}
\setlength{\headheight}{15pt}

\pagestyle{fancy}
\fancyhf{}
\fancyhead[L]{\textcolor{darkaccent}{Чек-лист по геометрии}}
\fancyhead[R]{\textcolor{darkaccent}{Трапеция}}
\fancyfoot[C]{\thepage}

\setlist[itemize]{leftmargin=1.8em}
\setlist[enumerate]{leftmargin=2em}

\newcommand{\ds}{\displaystyle}
\newcommand{\med}{m}
\newcommand{\Area}{S}
\newcommand{\half}{\frac12}

\tkzSetUpLine[line width=0.9pt,color=accent]
\tkzSetUpPoint[color=accent,fill=accent,size=3]
\tkzSetUpLabel[color=black,font=\small]

\begin{document}

\begin{titlepage}
\centering
\vspace*{2.2cm}

{\Huge\bfseries \textcolor{darkaccent}{Чек-лист}\par}
\vspace{0.8cm}
{\LARGE\bfseries \textcolor{textgray}{Трапеция: площадь, средняя линия,}\\[0.2em]
\textcolor{textgray}{равнобедренная трапеция, диагонали, окружность}\par}

\vspace{2cm}

{\Large \textcolor{textgray}{Лёвин Артём Александрович эксклюзивно для Ксении Клыковой}\par}

\vspace{0.8cm}
{\large \today\par}

\vfill

\begin{tikzpicture}[scale=1]
\draw[line width=1.2pt, color=accent]
(-6,0) .. controls (-4,0.9) and (-2,-0.5) .. (0,0.35)
.. controls (1.5,0.95) and (3,-0.2) .. (6,0.55);
\draw[line width=0.7pt, color=darkaccent]
(-5.8,-0.35) .. controls (-3.8,0.15) and (-1.8,-0.8) .. (0.1,-0.25)
.. controls (2.2,0.35) and (4.2,-0.45) .. (5.8,-0.05);
\end{tikzpicture}

\vspace*{0.8cm}
\end{titlepage}

\section*{\textcolor{darkaccent}{1. Основные определения и обозначения}}

\begin{itemize}
    \item \textbf{Трапеция} --- четырёхугольник, у которого две стороны параллельны.
    \item \textbf{Основания трапеции} --- параллельные стороны. Ниже обозначаем их как $AD$ и $BC$, где $AD \parallel BC$.
    \item \textbf{Боковые стороны} --- непараллельные стороны: $AB$ и $CD$.
    \item \textbf{Высота трапеции} --- расстояние между основаниями; это перпендикуляр, опущенный с одного основания на другое.
    \item \textbf{Средняя линия трапеции} --- отрезок, соединяющий середины боковых сторон.
    \item \textbf{Равнобедренная трапеция} --- трапеция, у которой боковые стороны равны: $AB=CD$.
\end{itemize}

\section*{\textcolor{darkaccent}{2. Базовые формулы, которые нужно знать наизусть}}

\subsection*{\textcolor{accent}{2.1. Площадь трапеции}}
\[
\Area_{\text{трап}}=\frac{AD+BC}{2}\cdot h=\med\cdot h,
\]
где
\[
\med=\frac{AD+BC}{2}
\]
--- средняя линия, а $h$ --- высота.

\subsection*{\textcolor{accent}{2.2. Средняя линия}}
\[
\med=\frac{AD+BC}{2}.
\]

Это очень удобная формула: если известна средняя линия и одно основание, то второе находится сразу.

\subsection*{\textcolor{accent}{2.3. Формула Герона для треугольника}}
Если у треугольника стороны $a$, $b$, $c$, а
\[
p=\frac{a+b+c}{2},
\]
то
\[
\Area_{\triangle}=\sqrt{p(p-a)(p-b)(p-c)}.
\]

Эта формула в теме трапеции часто нужна после \textbf{параллельного переноса}, когда задача сводится к треугольнику.

\section*{\textcolor{darkaccent}{3. Главная идея занятия: сводить трапецию к треугольнику}}

\subsection*{\textcolor{accent}{3.1. Параллельный перенос}}
Один из самых полезных приёмов в задачах на трапецию:

\begin{itemize}
    \item переносим \textbf{боковую сторону} или \textbf{диагональ} параллельно самой себе;
    \item получаем \textbf{треугольник}, у которого известны или легко находятся стороны;
    \item высоту трапеции выражаем как высоту этого треугольника;
    \item затем используем либо формулу Герона, либо свойства прямоугольного/равнобедренного треугольника.
\end{itemize}

\subsection*{\textcolor{accent}{3.2. Почему это работает}}
С треугольником почти всегда проще:
\begin{itemize}
    \item есть Герон;
    \item есть теорема Пифагора;
    \item есть свойства углов;
    \item есть медианы, высоты, биссектрисы;
    \item легко считать площадь двумя способами.
\end{itemize}

\subsection*{\textcolor{accent}{3.3. Метод равенства площадей}}
Если одна и та же площадь посчитана по-разному, это даёт уравнение.

Обычно:
\[
\Area_{\triangle}=\sqrt{p(p-a)(p-b)(p-c)}
\]
и одновременно
\[
\Area_{\triangle}=\frac12 \cdot a \cdot h.
\]

Отсюда
\[
h=\frac{2\Area_{\triangle}}{a}.
\]

Именно так часто ищут высоту трапеции после сведения к треугольнику.

\section*{\textcolor{darkaccent}{4. Равнобедренная трапеция: что обязательно помнить}}

\begin{itemize}
    \item Боковые стороны равны:
    \[
    AB=CD.
    \]
    \item Диагонали равны:
    \[
    AC=BD.
    \]
    \item Углы при каждом основании попарно равны.
    \item Если опустить высоты на большее основание, то оно распадается на три части:
    \[
    \text{(малый кусок)}+\text{(средний кусок)}+\text{(малый кусок)}.
    \]
\end{itemize}

\subsection*{\textcolor{accent}{4.1. Очень важный разрез основания}}
Пусть в равнобедренной трапеции $AD \parallel BC$, причём $AD>BC$. Опустим из $B$ высоту $BH$ на $AD$.

Тогда:
\[
HD=\frac{AD-BC}{2},
\qquad
AH=\frac{AD+BC}{2}=\med.
\]

Это нужно помнить: \textbf{в равнобедренной трапеции отрезок от левого конца большего основания до основания высоты равен средней линии}.

\subsection*{\textcolor{accent}{4.2. Полезный вывод}}
Если известны диагональ и угол с основанием, то часто выгодно рассматривать прямоугольный треугольник, в котором один катет равен средней линии:
\[
AH=\med.
\]

\section*{\textcolor{darkaccent}{5. Диагонали трапеции}}

\subsection*{\textcolor{accent}{5.1. Если известны диагонали и основания}}
Тогда полезно:
\begin{itemize}
    \item перенести диагональ параллельно самой себе;
    \item получить треугольник, у которого третья сторона выражается через сумму или разность оснований;
    \item далее искать высоту через площадь треугольника.
\end{itemize}

\subsection*{\textcolor{accent}{5.2. Если трапеция равнобедренная}}
Тогда автоматически:
\[
AC=BD.
\]
Это резко упрощает вычисления.

\subsection*{\textcolor{accent}{5.3. Если диагонали перпендикулярны}}
Это сильное условие. В равнобедренной трапеции после параллельного переноса можно получить \textbf{прямоугольный равнобедренный треугольник}.

Тогда:
\begin{itemize}
    \item высота к гипотенузе является одновременно медианой;
    \item медиана к гипотенузе равна её половине;
    \item если гипотенуза выражается через основания, высота находится очень быстро.
\end{itemize}

\section*{\textcolor{darkaccent}{6. Окружность и трапеция}}

В занятии обсуждалась трапеция, \textbf{описанная около окружности}.

\subsection*{\textcolor{accent}{6.1. Что значит ``описана около окружности''}}
Это значит, что окружность касается всех четырёх сторон трапеции.

\subsection*{\textcolor{accent}{6.2. Что важно заметить сразу}}
Если трапеция описана около окружности, то в рамках школьных задач очень часто рассматривается именно \textbf{равнобедренная трапеция}. Это естественная и наиболее удобная конфигурация.

\subsection*{\textcolor{accent}{6.3. Как искать высоту через радиус}}
Если окружность касается двух параллельных оснований, то расстояние между основаниями связано с радиусом.

В стандартной конфигурации, когда центр окружности расположен между основаниями,
\[
h=2r.
\]

Это мгновенный вывод из того, что радиусы, проведённые к обоим основаниям, перпендикулярны основаниям и вместе дают всю высоту.

\subsection*{\textcolor{accent}{6.4. Когда центр окружности нарисован ``не внутри''}}
Даже если на чертеже центр кажется смещённым относительно фигуры, принцип остаётся тем же:
\begin{itemize}
    \item радиус к касательной перпендикулярен касательной;
    \item высота выражается через суммы или разности перпендикулярных отрезков;
    \item задача снова сводится к прямоугольным треугольникам.
\end{itemize}

\section*{\textcolor{darkaccent}{7. Алгоритмы решения типовых задач}}

\subsection*{\textcolor{accent}{7.1. Если даны основания и боковые стороны, а нужно найти площадь}}
\begin{enumerate}
    \item Записать
    \[
    \Area=\frac{AD+BC}{2}\cdot h.
    \]
    \item Понять, что неизвестна только высота $h$.
    \item Сделать параллельный перенос боковой стороны.
    \item Получить треугольник с известными тремя сторонами.
    \item Найти его площадь по Герону.
    \item Приравнять к
    \[
    \frac12 \cdot \text{основание} \cdot h.
    \]
    \item Найти $h$, затем площадь трапеции.
\end{enumerate}

\subsection*{\textcolor{accent}{7.2. Если даны основания и диагонали, а нужно найти площадь}}
\begin{enumerate}
    \item Снова начать с формулы площади трапеции.
    \item Осознать, что нужна высота.
    \item Перенести диагональ параллельно самой себе.
    \item Получить треугольник, где две стороны --- диагонали, а третья --- сумма оснований.
    \item Найти высоту через площадь треугольника.
\end{enumerate}

\subsection*{\textcolor{accent}{7.3. Если трапеция равнобедренная и известен угол}}
\begin{enumerate}
    \item Опустить высоту.
    \item Использовать факт
    \[
    AH=\med.
    \]
    \item Рассмотреть прямоугольный треугольник.
    \item Применить тригонометрию или свойства треугольника $30^\circ$--$60^\circ$--$90^\circ$.
\end{enumerate}

\subsection*{\textcolor{accent}{7.4. Если есть окружность}}
\begin{enumerate}
    \item Сразу провести радиусы к точкам касания.
    \item Отметить прямые углы.
    \item Разбить высоту на удобные части.
    \item Свести задачу к одному или двум прямоугольным треугольникам.
\end{enumerate}

\section*{\textcolor{darkaccent}{8. Мини-примеры для быстрого повторения}}

\subsection*{\textcolor{accent}{Пример 1. Средняя линия}}
Если основания равны $14$ и $26$, то
\[
\med=\frac{14+26}{2}=20.
\]

\subsection*{\textcolor{accent}{Пример 2. Площадь}}
Если основания $8$ и $18$, а высота $6$, то
\[
\Area=\frac{8+18}{2}\cdot 6=13\cdot 6=78.
\]

\subsection*{\textcolor{accent}{Пример 3. Равнобедренная трапеция}}
Если основания равны $10$ и $24$, то после опускания высот:
\[
\frac{24-10}{2}=7
\]
--- каждый ``краевой'' кусок большего основания.

А отрезок от левого конца большего основания до основания высоты:
\[
\frac{24+10}{2}=17,
\]
то есть это средняя линия.

\subsection*{\textcolor{accent}{Пример 4. Описанная около окружности трапеция}}
Если радиус окружности $r=5$, то высота трапеции:
\[
h=2r=10.
\]

\section*{\textcolor{darkaccent}{9. Типичные ошибки}}

\begin{itemize}
    \item Путать \textbf{среднюю линию} и \textbf{высоту}.
    \item Писать
    \[
    \Area=(AD+BC)\cdot h
    \]
    вместо
    \[
    \Area=\frac{AD+BC}{2}\cdot h.
    \]
    \item В равнобедренной трапеции забывать, что
    \[
    AC=BD.
    \]
    \item После опускания высот ошибочно считать, что большее основание делится пополам. На самом деле пополам делится не оно, а конструкция разбирается через
    \[
    \frac{AD-BC}{2}
    \quad \text{и} \quad
    \frac{AD+BC}{2}.
    \]
    \item После параллельного переноса не подписывать равные и параллельные отрезки.
    \item В задачах с окружностью забывать, что радиус к касательной перпендикулярен касательной.
    \item Использовать теорему Пифагора не в прямоугольном треугольнике.
    \item В формуле Герона ошибаться в полупериметре:
    \[
    p=\frac{a+b+c}{2},
    \]
    а не $p=a+b+c$.
\end{itemize}

\section*{\textcolor{darkaccent}{10. Один опорный чертёж}}

\begin{center}
\begin{tikzpicture}[scale=0.9]
\tkzDefPoint(0,0){A}
\tkzDefPoint(8,0){D}
\tkzDefPoint(2,3){B}
\tkzDefPoint(6,3){C}
\tkzDefProjPoint(B)(A,D) \tkzGetPoint{H}
\tkzDefMidPoint(B,C) \tkzGetPoint{M}

\tkzDrawPolygon(A,B,C,D)
\tkzDrawSegment(B,H)
\tkzDrawSegment(A,C)
\tkzDrawSegment(B,D)
\tkzDrawSegment(M,H)

\tkzMarkRightAngle[size=0.25](B,H,D)

\tkzLabelPoints[below](A,D,H)
\tkzLabelPoints[above](B,C,M)

\node[below] at ($(A)!0.5!(D)$) {$AD$};
\node[above] at ($(B)!0.5!(C)$) {$BC$};
\node[left] at ($(B)!0.5!(H)$) {$h$};
\node[right] at ($(M)!0.5!(H)$) {$\med$};
\end{tikzpicture}
\end{center}

\vspace{-0.5em}
\begin{center}
\small
На этом рисунке: $AD \parallel BC$, $BH=h$, средняя линия соединяет середины боковых сторон.
\end{center}

\section*{\textcolor{darkaccent}{11. Краткая памятка перед решением задач}}

\begin{enumerate}
    \item Всегда начинайте с вопроса: \textbf{что именно неизвестно?}
    \item Если нужна площадь, почти наверняка нужна \textbf{высота}.
    \item Если высота не дана, ищите способ свести задачу к \textbf{треугольнику}.
    \item В равнобедренной трапеции сразу вспоминайте:
    \[
    AC=BD,\qquad AH=\frac{AD+BC}{2},\qquad HD=\frac{AD-BC}{2}.
    \]
    \item Если в условии есть окружность, сразу проводите \textbf{радиусы к точкам касания}.
    \item Если известны три стороны треугольника, не забывайте про \textbf{Герона}.
\end{enumerate}

\newpage

\section*{\centering \textcolor{darkaccent}{Тренировочный блок}}

\textbf{Указание.} Во всех задачах аккуратно делайте чертёж и явно подписывайте известные отрезки.

\subsection*{\textcolor{accent}{Простые задачи}}

\begin{enumerate}[label=\textbf{\arabic*.}]
    \item Основания трапеции равны $12$ и $20$, высота равна $7$. Найдите:
    \begin{enumerate}[label=\asbuk*)]
        \item среднюю линию;
        \item площадь трапеции.
    \end{enumerate}

    \item Основания равнобедренной трапеции равны $10$ и $24$. Найдите:
    \begin{enumerate}[label=\asbuk*)]
        \item среднюю линию;
        \item каждый из двух крайних отрезков большего основания после опускания высот.
    \end{enumerate}

    \item В равнобедренной трапеции диагональ равна $14$ и образует с основанием угол $30^\circ$. Найдите среднюю линию трапеции.
\end{enumerate}

\subsection*{\textcolor{accent}{Средний уровень}}

\begin{enumerate}[label=\textbf{\arabic*.},start=4]
    \item В трапеции основания равны $16$ и $44$, боковые стороны равны $17$ и $25$. Найдите площадь трапеции.

    \item В трапеции основания равны $4$ и $11$, диагонали равны $9$ и $12$. Найдите площадь трапеции.

    \item Трапеция описана около окружности радиуса $6$. Найдите высоту трапеции.
\end{enumerate}

\subsection*{\textcolor{accent}{Выше среднего}}

\begin{enumerate}[label=\textbf{\arabic*.},start=7]
    \item В равнобедренной трапеции основания равны $24$ и $40$, а диагонали взаимно перпендикулярны. Найдите:
    \begin{enumerate}[label=\asbuk*)]
        \item высоту трапеции;
        \item площадь трапеции.
    \end{enumerate}

    \item В равнобедренной трапеции проведена высота $BH$ к большему основанию $AD$. Известно, что $AH=13$, а $HD=5$. Найдите основания трапеции.

    \item В равнобедренной трапеции диагональ равна $10$ и образует с большим основанием угол $60^\circ$. Найдите:
    \begin{enumerate}[label=\asbuk*)]
        \item среднюю линию;
        \item площадь трапеции, если высота равна $5\sqrt{3}$.
    \end{enumerate}
\end{enumerate}

\subsection*{\textcolor{accent}{Сложная задача}}

\begin{enumerate}[label=\textbf{\arabic*.},start=10]
    \item В равнобедренной трапеции основания равны $14$ и $40$, а окружность радиуса $25$ касается всех её сторон. Докажите, что высота трапеции выражается через радиус, и найдите эту высоту.
\end{enumerate}

\newpage

\section*{\centering \textcolor{darkaccent}{Ответы}}

\renewcommand{\arraystretch}{1.35}

\begin{table}[h!]
\centering
\caption{Краткие ответы к тренировочному блоку}
\begin{tabularx}{\linewidth}{>{\bfseries}c X}
\toprule
№ & Ответ \\ 
\midrule
1 & \textbf{а)} $\med=16$; \textbf{б)} $\Area=112$. \\
2 & \textbf{а)} $\med=17$; \textbf{б)} по $7$. \\
3 & $\med=7\cos 30^\circ=7\cdot \dfrac{\sqrt3}{2}= \dfrac{7\sqrt3}{2}$. \\
4 & $\Area=336$. \\
5 & $\Area=54$. \\
6 & $h=12$. \\
7 & \textbf{а)} $h=32$; \textbf{б)} $\Area=1024$. \\
8 & $BC=AH-HD=8$, \quad $AD=AH+HD=18$. \\
9 & \textbf{а)} $\med=5$; \textbf{б)} $\Area=25\sqrt3$. \\
10 & $h=2r=50$. \\
\bottomrule
\end{tabularx}
\end{table}

\vfill

{\small
\textcolor{softgray}{Памятка: если решение ``не идёт'', ищите треугольник, а в нём --- прямой угол, равные стороны, известные углы или возможность применить формулу Герона.}
}

\end{document}